% SOURCE_AUTHORITY: SGA5 (1).pdf, PDF pages 3--4 (Roman pages III--IV; scan-only authority).
% SOURCE_SHA256: B256EBD072A8C68209518412A263C9289C6F1854A346733D86F885930D5FE6CA
\thispagestyle{empty}
\section*{\uline{Introducción}}

Este volumen reúne las exposiciones del seminario de Grothendieck celebrado en el IHES en 1965/66, distribuidas inicialmente en forma de notas mimeografiadas. Con respecto a la versión primitiva, los únicos cambios importantes conciernen a la exposición II, que no se reproduce, y a la exposición III, que ha sido reescrita por completo y ampliada con un apéndice numerado III~B. Salvo algunas modificaciones de detalle y adiciones de notas a pie de página, las demás exposiciones se han dejado tal como estaban.

El núcleo del seminario está constituido por la fórmula de Lefschetz en cohomología étale (III, III~B, XII) y por su aplicación a la interpretación cohomológica de las funciones $L$ (XIV). Todos los resultados anunciados por Grothendieck en su exposición ante el seminario Bourbaki~[2] se demuestran aquí por completo. Las fórmulas de las trazas establecidas en (III, III~B y XII), por vías diferentes, son más generales de lo necesario para demostrar únicamente la racionalidad de las funciones $L$. En la exposición de Deligne (SGA~$4^{1/2}$ Rapport) figura una demostración completa de esta última, considerablemente más breve, en la que se sigue el método de Grothendieck de las exposiciones XII y XIV, pero despojado de toda generalidad superflua. Esperamos, no obstante, que las fórmulas de III y III~B puedan servir en otras situaciones. En cuanto al resto del seminario, consta de dos exposiciones sobre la teoría del paso al límite que da origen a la cohomología $\ell$-ádica (V, VI), y de diversos complementos al formalismo de dualidad (I, VII) y a la fórmula de Lefschetz (VII, X).

He aquí, con mayor precisión, el contenido de este volumen.

La exposición I es independiente del resto del seminario. Su resultado principal afirma que, sobre un esquema regular que satisface ciertas hipótesis adicionales de naturaleza local, y siempre que se disponga de la resolución de singularidades y del teorema de pureza, el haz constante de valor $\mathbb Z/n\mathbb Z$, para $n$ primo con las características residuales, es dualizante (I.3.4.1). Este teorema vale, en particular, para los esquemas regulares excelentes de característica nula, gracias a Hironaka y a los resultados de Artin (SGA~4.XIX). El caso de un esquema regular de dimensión $1$ se

% SOURCE_PAGE_BOUNDARY: the printed sentence continues from Roman p. III to p. IV here.
\clearpage
\noindent trata por separado, mediante otro método, muy elemental (véase también (SGA~$4^{1/2}$ Dualité)). Señalemos que, mediante un argumento diferente que no utiliza ni resolución ni pureza, Deligne demuestra en (SGA~$4^{1/2}$ Th. Finitude) que, si $f:X\longrightarrow S$ presenta a $X$ como un esquema de tipo finito sobre un esquema regular $S$ de dimensión $0$ o $1$, el complejo $f^!(\mathbb Z/n\mathbb Z)_S$ (con $n$ primo con las características residuales) es dualizante. Se ignora, sin embargo, si este resultado, más útil que el teorema de bidualidad de la exposición I, con el cual coincide parcialmente sin contenerlo, se extiende al caso de un esquema de base $S$ excelente, regular y de dimensión $>1$.

Como se ha dicho más arriba, la exposición II, que llevaba por título ``Fórmulas de Künneth para la cohomología con soportes cualesquiera'', no figura en este volumen. Redactada por L.~Illusie a partir de notas manuscritas de Grothendieck, estaba dedicada a teoremas genéricos sobre la propiedad de ser cohomológicamente propio y sobre la aciclicidad local, pero estos solo se demostraban bajo hipótesis de resolución. Una demostración de los mismos resultados sin esas hipótesis, obtenida posteriormente por Deligne (SGA~$4^{1/2}$ Th. Finitude), hizo innecesaria la publicación de esta exposición. Algunos complementos, relativos, por ejemplo, al caso de los divisores de cruzamientos normales, se incorporaron a un apéndice de (loc.~cit.). Añadamos que en la exposición de la Sra.~Raynaud (SGA~1 XIII) se demuestran, para el caso de los haces de conjuntos y de grupos no conmutativos, teoremas análogos a los de (loc.~cit.).

La exposición III contiene una demostración de la fórmula de Lefschetz--Verdier para las correspondencias cohomológicas entre complejos de haces sobre esquemas de tipo finito sobre un cuerpo; en la exposición III~B se muestra cómo calcular los términos locales de la fórmula en ciertos casos, por ejemplo el de la correspondencia de Frobenius sobre las curvas. Remitimos al lector a las introducciones de las exposiciones III y III~B para más detalles sobre su contenido.

La exposición IV, de A.~Grothendieck, dedicada a la construcción de la clase de cohomología asociada a un ciclo, fue redactada por P.~Deligne y publicada en SGA~$4^{1/2}$.
